% ============================================================
% Chương 3. Mô tả bài toán và dữ liệu
% Nguồn: docs/overview_and_demo.md mục 1, docs/report_tv1.md mục 3.3,
%        file đặc tả 8 BR (OULAD – Giai đoạn 2 Đặc tả 8 BR.docx)
% ============================================================
\section{Mô tả project}
\label{sec:mo-ta}

\subsection{Bộ dữ liệu OULAD}
\label{sec:nguon-du-lieu}
Hệ thống dùng bộ dữ liệu Open University Learning Analytics Dataset (OULAD) do Open
University, một trường đại học đào tạo từ xa ở Anh, công bố năm 2017 theo giấy phép
CC BY 4.0~\cite{kuzilek2017oulad}. Bộ dữ liệu ghi lại quá trình học của 28.785 sinh viên
trong 7 môn học (ký hiệu AAA đến GGG), mỗi môn mở ở một hoặc nhiều kỳ trong hai năm 2013 và
2014, tổng cộng 22 đợt học. Mã kỳ kết thúc bằng B là kỳ bắt đầu vào tháng 2, bằng J là kỳ
bắt đầu vào tháng 10. Thông tin cá nhân của sinh viên đã được ẩn danh.

Dữ liệu gồm 7 file CSV, mô tả ở Bảng~\ref{tab:file-nguon}. Ba file đầu mô tả môn học, bài
kiểm tra và học liệu; bốn file sau ghi lại sinh viên và hoạt động của họ. File lớn nhất là
\texttt{studentVle}, ghi số click của từng sinh viên trên từng học liệu theo ngày, với hơn
10,6 triệu dòng và tổng cộng 39.605.099 click.

\begin{table}[htbp]
\centering
\small
\begin{tabularx}{\textwidth}{|l|X|r|}
\hline
\textbf{File} & \textbf{Nội dung, một dòng là} & \textbf{Số dòng} \\ \hline
\texttt{courses} & Một đợt học: mã môn, mã kỳ, số ngày của đợt & 22 \\ \hline
\texttt{assessments} & Một bài kiểm tra của một đợt: loại bài (TMA, CMA, Exam), hạn nộp, trọng số & 206 \\ \hline
\texttt{vle} & Một học liệu trên VLE của một đợt: loại học liệu (trang nội dung, diễn đàn, quiz\ldots) & 6.364 \\ \hline
\texttt{studentInfo} & Một sinh viên trong một đợt: đặc điểm cá nhân và kết quả cuối kỳ & 32.593 \\ \hline
\texttt{studentRegistration} & Một lượt đăng ký: ngày đăng ký, ngày rút (nếu có) & 32.593 \\ \hline
\texttt{studentAssessment} & Một lần nộp bài: ngày nộp, điểm, có phải điểm chuyển từ đợt trước không & 173.912 \\ \hline
\texttt{studentVle} & Số click của một sinh viên trên một học liệu trong một ngày & 10.655.280 \\ \hline
\end{tabularx}
\caption{Bảy file CSV của bộ dữ liệu OULAD}
\label{tab:file-nguon}
\end{table}

Đây là bản xuất tĩnh, không thay đổi sau khi công bố, nên nhóm không dựng một cơ sở dữ liệu
vận hành (OLTP) để mô phỏng nguồn mà đọc thẳng các file CSV vào pipeline. Nếu nguồn là một
hệ thống đang chạy và dữ liệu thay đổi liên tục, kho sẽ cần nạp tăng dần theo các thay đổi
(CDC); hướng này được bàn ở mục~\ref{sec:huong-phat-trien}. Dù cùng thuộc một bộ dữ liệu,
7 file vẫn khác nhau về khóa, mức độ sạch và cách ghi ngày, nên việc hợp nhất chúng vào kho
cũng gặp các vấn đề như khi tích hợp nhiều nguồn (mục~\ref{sec:van-de-du-lieu}).

\subsection{Yêu cầu nghiệp vụ}
\label{sec:br}
% Nguồn: OULAD – Giai đoạn 2 Đặc tả 8 BR.docx

\subsubsection{Thuật ngữ và quy ước chung}
Các yêu cầu nghiệp vụ dùng chung những thuật ngữ và quy ước sau:
\begin{itemize}
  \item \textbf{Đợt học}: một môn mở trong một kỳ, ví dụ BBB 2014J. Bộ dữ liệu có 22 đợt.
        Mọi tỷ lệ được tính trong từng đợt rồi mới đem so sánh.
  \item \textbf{Lượt đăng ký}: một sinh viên học một môn trong một đợt. Một sinh viên
        có thể có nhiều lượt nếu học nhiều môn hoặc học lại.
  \item \textbf{Ngày và tuần}: ngày đếm từ ngày khai giảng (ngày 0), ngày âm là trước
        khai giảng. Tuần 0 gồm ngày 0--6, tuần 1 gồm ngày 7--13, sớm nhất là tuần $-4$.
  \item \textbf{TMA và CMA}: bài tập do giảng viên chấm (Tutor Marked Assessment) và bài
        trắc nghiệm do máy chấm (Computer Marked Assessment). Bài thi cuối kỳ (Exam)
        không được dùng vì chỉ 6/22 đợt có điểm thi.
  \item \textbf{VLE}: hệ thống học trực tuyến (Virtual Learning Environment). Mỗi lần
        sinh viên bấm vào một học liệu được ghi là một click.
  \item \textbf{IMD}: chỉ số mức độ khó khăn của khu vực sinh viên sinh sống (Index of
        Multiple Deprivation), chia thành 10 nhóm từ 0--10\% (khó khăn nhất) đến 90--100\%.
  \item \textbf{Đạt}: đạt môn là có kết quả Pass hoặc Distinction; đạt một bài là từ 40 điểm.
  \item \textbf{Còn học}: sinh viên chưa rút khỏi môn tính đến thời điểm đang xét. Sinh viên
        rút trước khai giảng không có dữ liệu học tập nên không được tính ở các chỉ số
        về hành vi học tập (BR2, BR3, BR4, BR7, BR8).
\end{itemize}
Các kết quả chỉ mô tả mối liên hệ trong dữ liệu. Báo cáo và dashboard không dùng câu
hàm ý nhân quả như ``click nhiều giúp điểm cao''.

\subsubsection{Tám yêu cầu nghiệp vụ}
\begin{table}[htbp]
\centering
\small
\begin{tabularx}{\textwidth}{|l|l|X|}
\hline
\textbf{Mã} & \textbf{Người dùng} & \textbf{Câu hỏi nghiệp vụ} \\ \hline
BR1 & Giảng viên & Mỗi bài tập có bao nhiêu sinh viên nộp, bao nhiêu nộp trễ, và điểm ra sao? \\ \hline
BR2 & Giảng viên & Sinh viên dùng loại học liệu nào nhiều nhất, và vào giai đoạn nào của đợt họ hoạt động nhiều hay ít? \\ \hline
BR3 & Giảng viên & Sinh viên click nhiều trước hạn bài TMA đầu tiên thì có hay được điểm bài đó cao hơn không? \\ \hline
BR4 & Cố vấn học tập & Tính đến cuối tuần này, sinh viên nào đang có dấu hiệu nguy cơ rớt hoặc bỏ môn? \\ \hline
BR5 & Cố vấn học tập & Sinh viên rút khỏi môn vào thời điểm nào, và đến mỗi mốc bao nhiêu phần trăm đã rút? \\ \hline
BR6 & Lãnh đạo khoa & Tỷ lệ đạt của nhóm IMD khó khăn nhất thấp hơn nhóm ít khó khăn nhất bao nhiêu, và thay đổi thế nào qua các kỳ? \\ \hline
BR7 & Sinh viên & Tính đến tuần này, mình đang ở đâu so với các bạn cùng đợt về hoạt động, số bài đã nộp và điểm? \\ \hline
BR8 & Sinh viên & Trong tuần này, mình dùng từng loại học liệu nhiều hay ít hơn so với các bạn đã đạt môn? \\ \hline
\end{tabularx}
\caption{Tám yêu cầu nghiệp vụ}
\label{tab:br}
\end{table}

Chi tiết từng yêu cầu:
\begin{itemize}
  \item \textbf{BR1 -- Tình hình nộp bài và điểm từng bài.} Với mỗi bài, hệ thống tính tỷ lệ nộp (số sinh viên đã nộp chia cho số sinh viên còn học
        đến hạn nộp), tỷ lệ nộp trễ (chỉ tính cho TMA) và phân bố điểm (trung vị, khoảng điểm
        của 50\% sinh viên ở giữa, phần trăm bài đạt). CMA không có tỷ lệ trễ vì ngày ghi của CMA
        không phải hạn nộp thật. Bài được chuyển điểm từ đợt trước không được tính.
  \item \textbf{BR2 -- Hoạt động trên VLE theo loại học liệu và theo tuần.} Số đo chính là số click trung bình mỗi sinh viên trong từng tuần: tổng click chia cho số
        sinh viên còn học đến cuối tuần, kể cả người không click. Kèm theo là tỷ trọng click của
        từng loại học liệu. Để so hai đợt có độ dài khác nhau (234--269 ngày), trục thời gian
        được đổi sang phần trăm độ dài đợt.
  \item \textbf{BR3 -- Hoạt động trước bài TMA đầu tiên và điểm bài đó.} Với mỗi sinh viên, hệ thống lấy tổng click trước hạn của bài TMA đầu tiên có trọng số và
        điểm của bài đó, rồi tính hệ số tương quan hạng Spearman ($\rho$) riêng cho từng đợt, kèm
        khoảng tin cậy 95\%. Chỉ có 19/22 đợt vì môn GGG không có TMA nào có trọng số. Đây là
        tương quan, không phải quan hệ nhân quả.
  \item \textbf{BR4 -- Danh sách sinh viên có nguy cơ, cập nhật theo tuần.} Mỗi tuần, sinh viên còn học được xét 4 dấu hiệu: tổng click thuộc nhóm 25\% thấp nhất của
        đợt; không click trong 2 tuần gần nhất; có TMA đã qua hạn mà chưa nộp; điểm trung bình
        dưới 40. Sinh viên có từ 2 dấu hiệu trở lên vào danh sách. Danh sách tuần $w$ chỉ dùng dữ
        liệu đến hết tuần $w$ và không dùng kết quả cuối kỳ. Kết quả cuối kỳ chỉ dùng sau khi đợt
        kết thúc, để đo độ chính xác và độ bao phủ của quy tắc.
  \item \textbf{BR5 -- Đường cong rút lui.} Hệ thống vẽ phần trăm lượt đăng ký đã rút theo từng ngày, từ trước khai giảng đến hết đợt,
        kèm ngày rút trung vị và tỷ lệ rút trước khai giảng. ``Rút'' ở đây là rút khỏi một môn,
        không phải bỏ học cả trường.
  \item \textbf{BR6 -- Khoảng cách tỷ lệ đạt giữa các nhóm IMD.} Hệ thống so tỷ lệ đạt và tỷ lệ Distinction của nhóm IMD 0--10\% với nhóm 90--100\%, tính
        khoảng cách bằng điểm phần trăm qua 4 kỳ. Chỉ dùng 3 môn có đủ 4 kỳ (BBB, DDD, FFF) để xu
        hướng không bị lệch do mỗi kỳ mở môn khác nhau. Kết quả chỉ báo ở mức nhóm.
  \item \textbf{BR7 -- Tiến độ tuần của sinh viên so với các bạn cùng đợt.} Sinh viên xem 4 chỉ số của mình tính đến tuần đang chọn (click trong tuần, tổng click,
        số bài đã nộp, điểm trung bình) và vị trí của mình so với phân vị 25\%, 50\%, 75\% của các
        bạn cùng đợt còn học. Nhóm dưới 30 sinh viên thì không hiện phân vị, để tránh đoán ra số
        của người khác.
  \item \textbf{BR8 -- So hoạt động của mình với nhóm đã đạt.} Với mỗi loại học liệu, sinh viên thấy số click của mình trong tuần đặt trên dải 25\%--75\%
        click của các bạn đã đạt môn. Trên OULAD, nhóm đạt lấy cùng đợt vì mọi đợt đã kết thúc;
        khi dùng thật thì lấy đợt trước của cùng môn, vì đợt đang học chưa có kết quả.
\end{itemize}

Đi kèm 8 yêu cầu là một bộ 20 câu hỏi mẫu, mỗi yêu cầu có ít nhất một câu, ví dụ ``tỷ lệ
nộp trễ chung của mọi bài TMA là bao nhiêu''. Đáp án của từng câu được tính thẳng từ 7 file
CSV gốc, không đi qua kho, nên bộ câu hỏi này dùng được để kiểm tra kho trả lời đúng. Kết quả
kiểm tra được trình bày ở chương~\ref{sec:danh-gia}.

\subsection{Các vấn đề của dữ liệu nguồn}
\label{sec:van-de-du-lieu}
Dữ liệu OULAD có một số vấn đề khiến nó không thể nạp thẳng vào kho; cách xử lý được
trình bày ở mục~\ref{sec:etl}.

File click \texttt{studentVle} không có khóa: cùng một sinh viên, học liệu và ngày có thể
nằm trên nhiều dòng, tổng cộng 3.810.465 trên 10.655.280 dòng. Giá trị thiếu được ghi bằng
dấu \texttt{"?"}, và cột nhóm IMD vừa thiếu ở 971 sinh viên vừa ghi một nhóm theo hai cách
(\texttt{10-20} và \texttt{10-20\%}).

Dữ liệu không có ngày lịch thật: mọi mốc thời gian là số ngày tính từ khai giảng, có thể
âm. Vì vậy cách quy ngày về tuần phải thống nhất trên mọi engine, nếu không cùng một truy
vấn sẽ ra kết quả khác nhau.

Một số dòng thiếu thông tin hoặc mâu thuẫn: 11 bài thi cuối kỳ không có hạn nộp, và 102
lượt đăng ký có kết quả cuối kỳ không khớp với ngày rút.

Riêng file \texttt{studentInfo} có vấn đề về cấu trúc bảng chứ không phải về giá trị. Mỗi
dòng là một sinh viên trong một đợt, nhưng file chứa cả đặc điểm của sinh viên (giới tính,
vùng, trình độ, nhóm IMD) lẫn thông tin của lượt đăng ký (nhóm tuổi, số lần học lại, kết
quả cuối kỳ). Đặc điểm sinh viên chỉ phụ thuộc vào mã sinh viên, không phụ thuộc vào đợt,
nên bị lặp lại ở mỗi đợt sinh viên đó học; bảng vi phạm dạng chuẩn 2 (2NF). Khi tách bảng
cũng phải xét kỹ từng cột: nhóm tuổi trông giống đặc điểm sinh viên nhưng thay đổi giữa các
lần đăng ký ở 72 sinh viên, nên thuộc về lượt đăng ký. Việc tách bảng theo phụ thuộc hàm
được trình bày ở mục~\ref{sec:chuan-hoa}.
